\chapter{Notation}
\markboth{Notation}{Notation}

This book covers eight fields that grew up apart, and their alphabets collide.
The tables below list every symbol that appears in more than one place, say what
it means, and name the chapter that introduces it. Where a letter genuinely
carries two meanings---$\mat{Q}$ and $\mat{R}$ are the clearest case---the table
says so instead of pretending otherwise, and the chapter concerned repeats the
warning at the point of use.

\section*{Conventions}

\textbf{Fonts.} Scalars are italic ($v$, $t$, $c$, $\lambda$), vectors are bold
lower case ($\pos$, $\vel$, $\state$, $\ctrl$, $\meas$) and matrices are bold
upper case ($\mat{A}$, $\mat{P}$, $\mat{\Sigma}$). The zero vector is
$\vect{0}$, the identity matrix $\mat{I}$ and the vector of ones $\vect{1}$;
their sizes follow from the context. Sets and spaces are calligraphic
($\Cspace$, $\Cfree$, $\mathcal{C}_i$, $\mathcal{W}$, $\Normal$), with the one
classical exception that the vertex and edge sets of a graph keep the plain
letters $V$ and $E$. Algorithm data structures are small capitals (\Open,
\Closed, \Focal), and names taken from code are typeset as \code{this}.

\textbf{Operators.} $\mat{A}\T$ is the transpose and $\mat{A}\inv$ the inverse;
$\norm{\vect{x}}$ is the Euclidean norm unless a subscript says otherwise, and
$\abs{\cdot}$ is either an absolute value or, for a finite set, its number of
elements ($\abs{V}$, $\abs{E}$). $\mat{Q}\succeq\mat{0}$ means that $\mat{Q}$ is
symmetric positive semidefinite, $\succ$ that it is positive definite.

\textbf{Subscripts and superscripts.} A subscript $k$ or $t$ is a time index:
$\state_k$ is the state at step $k$, that is at the instant $k\,\dt$, and
$\pi[t]$ is the vertex occupied at step $t$. A subscript $i$ or $j$ is an agent
index: $\pi_i$ is the path of agent $a_i$, $\pos_i$ its position, $r_i$ its
radius. When both are needed the agent moves up and the time stays down, as in
$\pos^i_k$, the position of agent $i$ at step $k$ of a horizon; the superscript
is an index, never an exponent, and a power is bracketed ($\mat{A}^{k}$ is a
power, $\pos^{i}$ is agent $i$). A star marks an optimal quantity
($\gcost^*$, $\hcost^*$, $C^*$, $J^*$, $t^*$). A hat marks an estimate produced
by a filter or a predictor ($\hat{\state}_k$, $\hat{\pos}_{T+h}$), a bar an
average ($\bar{\state}$, $\bar{x}$), and a superscript minus a \emph{prediction}
made before the measurement of the current step has been used
($\hat{\state}_k^-$, $\mat{P}_k^-$). Parenthesised superscripts count samples:
$\state^{(i)}$ is the $i$-th particle or sigma point.

\textbf{Time.} Search-based chapters (Parts~II and~III) count time in unit
steps, $t=0,1,2,\dots$, and one action takes one step. Control and estimation
chapters (Parts~IV, VI and~VII) count steps of the sampling interval $\dt$
seconds, $k=0,1,2,\dots$, so step $k$ is the instant $k\,\dt$. The two views
meet in \cref{ch:ch02}.

\section*{General mathematics}

\begin{center}\small
\begin{tabularx}{\textwidth}{@{}lYl@{}}
\toprule
\textbf{Symbol} & \textbf{Meaning} & \textbf{Where introduced}\\
\midrule
$\R$, $\N$, $\Z$ & real numbers, natural numbers including $0$, integers; $\R^n$ is the space of $n$-vectors & \cref{ch:ch02}\\
$\vect{x}$, $\mat{A}$ & vector (bold lower case), matrix (bold upper case) & \cref{ch:ch02}\\
$\vect{0}$, $\vect{1}$, $\mat{I}$ & zero vector, vector of ones, identity matrix & \cref{ch:ch02}\\
$\mat{A}\T$, $\mat{A}\inv$ & transpose, inverse & \cref{ch:ch02}\\
$\norm{\vect{x}}$, $\abs{a}$ & Euclidean norm; absolute value, or cardinality of a set & \cref{ch:ch02}\\
$\vect{a}\cdot\vect{b}$, $\vect{a}\times\vect{b}$ & dot product; cross product (a signed scalar in the plane) & \cref{ch:ch02}\\
$\vect{a}^{\perp}$ & $\vect{a}=(a_x,a_y)$ rotated a quarter turn counter-clockwise, $(-a_y,a_x)$ & \cref{ch:ch02}\\
$A\oplus B$ & Minkowski sum $\set{\vect{a}+\vect{b} : \vect{a}\in A,\ \vect{b}\in B}$; obstacle inflation & \cref{ch:ch02}\\
$\set{\cdot}$, $\mathcal{S}$ & a set; calligraphic letters name sets and spaces & \cref{ch:ch02}\\
$\argmin$, $\argmax$ & the argument at which a function attains its minimum or maximum & \cref{ch:ch02}\\
$\diag(\cdot)$, $\tr(\cdot)$ & diagonal matrix built from a list; trace of a matrix & \cref{ch:ch02}\\
$\bigO{f(n)}$ & asymptotic upper bound as the input size grows & \cref{ch:ch02}\\
$\hat{x}$, $\bar{x}$, $x^*$ & estimate of $x$; average of $x$; optimal value of $x$ & \cref{ch:ch04,ch:ch18}\\
$x^-$ & value predicted before the current measurement is used & \cref{ch:ch18}\\
\bottomrule
\end{tabularx}
\end{center}

\section*{Graphs, grids and search}

\begin{center}\small
\begin{tabularx}{\textwidth}{@{}lYl@{}}
\toprule
\textbf{Symbol} & \textbf{Meaning} & \textbf{Where introduced}\\
\midrule
$G=(V,E)$ & graph with vertex set $V$ and edge set $E\subseteq V\times V$ & \cref{ch:ch02}\\
$c(u,v)$ & cost (weight) of the edge from $u$ to $v$; costs are non-negative & \cref{ch:ch02}\\
$\Succ(u)$, $\Pred(u)$ & successors and predecessors of $u$; $b$ is the largest out-degree & \cref{ch:ch02}\\
$G_4$, $G_8$ & 4- and 8-connected grid graph; a diagonal move costs $\sqrt2$ & \cref{ch:ch02}\\
$s$, $g$ & start vertex and goal vertex; \cref{ch:ch03} writes the target $t$, \cref{ch:ch04,ch:ch09} write the goal $\gamma$ & \cref{ch:ch02}\\
$\pi=(v_0,\dots,v_k)$ & a path; $\cost(\pi)$ its cost, $L(\pi)$ its Euclidean length & \cref{ch:ch02}\\
$\dist(u,v)$ & shortest-path distance; \cref{ch:ch03} writes $\delta(s,v)$ & \cref{ch:ch02}\\
$(v,t)$, $T$ & space-time state (vertex $v$ at step $t$); planning horizon $T$ & \cref{ch:ch02}\\
$\gcost(n)$, $\gcost^*(n)$ & cost of the best path from $s$ to node $n$ found so far; its optimum & \cref{ch:ch04}\\
$\hcost(n)$, $\hcost^*(n)$ & heuristic estimate of the cost to go; the true cost to go & \cref{ch:ch04}\\
$\fcost(n)$ & priority of a node, $\fcost=\gcost+\hcost$ & \cref{ch:ch04}\\
$C^*=\hcost^*(s)$ & cost of an optimal solution & \cref{ch:ch04}\\
\Open, \Closed & the frontier (priority queue) and the set of expanded nodes & \cref{ch:ch03}\\
\Focal & the band $\set{n\in\Open : \fcost(n)\le w\,\fcost_{\min}}$ of a focal search & \cref{ch:ch10}\\
$\fcost_{\min}$, $\mathrm{LB}$ & smallest $\fcost$ in \Open; the lower bound it provides on $C^*$ & \cref{ch:ch10}\\
$w\ge1$ & suboptimality factor of weighted \astar and of focal search & \cref{ch:ch04,ch:ch10}\\
$\eps$ & inflation factor of \arastar; $w=1+\eps$ in $\astar_\eps$ & \cref{ch:ch06,ch:ch10}\\
$\rhs(s)$ & one-step look-ahead value; $\gcost(s)=\rhs(s)$ is local consistency & \cref{ch:ch05}\\
$\key(s)$, $k_m$ & two-component priority of \lpastar/\dstarlite; the key modifier & \cref{ch:ch05}\\
\bottomrule
\end{tabularx}
\end{center}

\section*{Multi-agent path finding}

\begin{center}\small
\begin{tabularx}{\textwidth}{@{}lYl@{}}
\toprule
\textbf{Symbol} & \textbf{Meaning} & \textbf{Where introduced}\\
\midrule
$k$, $a_i$ & number of agents; agent $i$, with start $s_i$ and goal $g_i$ (or $\gamma_i$) & \cref{ch:ch07}\\
$\pi_i$, $\pi_i[t]$ & time-indexed path of $a_i$; the vertex it occupies at step $t$ & \cref{ch:ch07}\\
$\cost(\pi_i)$ & cost of a path: the time of its final arrival at the goal & \cref{ch:ch07}\\
$T_i$, $T$ & last stored index of $\pi_i$; horizon $T=\max_i T_i$ of the plan & \cref{ch:ch07}\\
$\Pi=(\pi_1,\dots,\pi_k)$ & a plan: one time-indexed path per agent & \cref{ch:ch07}\\
$\sumcost(\Pi)$ & sum of costs, $\sum_i\cost(\pi_i)$ & \cref{ch:ch07}\\
$\makespan(\Pi)$ & makespan, $\max_i\cost(\pi_i)$ & \cref{ch:ch07}\\
$\langle a_i,a_j,v,t\rangle$ & vertex conflict: both agents occupy $v$ at time $t$ & \cref{ch:ch07}\\
$\langle a_i,a_j,u,v,t\rangle$ & edge or swapping conflict along the edge $\set{u,v}$ between $t$ and $t+1$ & \cref{ch:ch07}\\
$\langle a_i,v,t\rangle$ & vertex constraint: $a_i$ may not occupy $v$ at time $t$ & \cref{ch:ch09}\\
$\mathcal{C}_i$ & the set of constraints imposed on agent $a_i$ & \cref{ch:ch09}\\
$N$ & node of the constraint tree, with $N.\mathcal{C}$, $N.\pi$ and $N.\cost$ & \cref{ch:ch09}\\
$\sigma$ & priority order: a permutation of the agents, $\sigma(1)$ first & \cref{ch:ch08}\\
$h_c$ & number of remaining conflicts, the \Focal heuristic of \ecbs & \cref{ch:ch10}\\
$C(v)$, $\Psi(v,v')$ & collision set of a joint vertex and of a joint move (\mstar) & \cref{ch:ch11}\\
$\phi_i$ & individual policy of agent $a_i$: its next vertex on a shortest path & \cref{ch:ch11}\\
\bottomrule
\end{tabularx}
\end{center}

\section*{Geometry and motion}

\begin{center}\small
\begin{tabularx}{\textwidth}{@{}lYl@{}}
\toprule
\textbf{Symbol} & \textbf{Meaning} & \textbf{Where introduced}\\
\midrule
$\mathcal{W}$, $\mathcal{O}$ & workspace and the obstacle region inside it & \cref{ch:ch02}\\
$\Cspace$, $\Cfree$, $\Cobs$ & configuration space, free space, obstacle region of the configuration space & \cref{ch:ch02}\\
$\disc{\vect{c}}{r}$ & closed disc (in 3D, ball) of radius $r$ centred at $\vect{c}$ & \cref{ch:ch02}\\
$r_i$, $R=r_A+r_B$ & radius of an agent; collision distance of a pair of agents & \cref{ch:ch02}\\
$\pos$, $\vel$, $\acc$ & position, velocity, acceleration & \cref{ch:ch02}\\
$v_{\max}$, $a_{\max}$ & speed limit and acceleration limit of a drone & \cref{ch:ch02}\\
$\dt$ & time step (sampling interval) in seconds & \cref{ch:ch02}\\
$\state$, $\ctrl$ & state and control input of a motion model & \cref{ch:ch02}\\
$\mat{A}$, $\mat{B}$ & state-transition and input matrix of the double integrator & \cref{ch:ch02}\\
$d_{ij}(t)$, $d_{\min}$ & separation of two agents at time $t$; its minimum over the execution & \cref{ch:ch02}\\
$d_{\mathrm{stop}}(v)$ & stopping distance $v^2/(2a_{\max})$ of a braking double integrator & \cref{ch:ch02}\\
$t^*$, $t_c$ & time of closest approach; time to collision of two constant-velocity objects & \cref{ch:ch02}\\
$\ttc$ & time horizon of the reactive layer (\cref{ch:ch01} writes $\ttc$ for the time to collision itself and $\ttc_{\mathrm{safe}}$ for the horizon) & \cref{ch:ch02,ch:ch12}\\
$\pos_{\mathrm{rel}}$, $\vel_{\mathrm{rel}}$ & relative position and velocity of a pair of agents & \cref{ch:ch02}\\
$CC_{A|B}$ & collision cone: relative velocities of $A$ that eventually hit $B$ & \cref{ch:ch12}\\
$\VO_{A|B}(\vel_B)$, $\VO^{\ttc}_{A|B}$ & velocity obstacle and its truncation at the horizon $\ttc$ & \cref{ch:ch12}\\
$\RVO_{A|B}(\vel_B,\vel_A)$ & reciprocal velocity obstacle: the cone with its apex at $\tfrac12(\vel_A+\vel_B)$ & \cref{ch:ch13}\\
$\ORCA^{\ttc}_{A|B}$ & \orca half-plane of velocities that $A$ may choose & \cref{ch:ch13}\\
$\vect{u}$, $\vect{n}$ & smallest change of the relative velocity that reaches the boundary of $\VO^{\ttc}$, and the outward unit normal there & \cref{ch:ch13}\\
$\vel^{\mathrm{pref}}$, $\vel^{\mathrm{new}}$ & preferred velocity and the velocity finally commanded & \cref{ch:ch13}\\
$V_s$, $V_a$, $V_d$, $V_r$ & possible, admissible, dynamic-window and resulting velocity sets & \cref{ch:ch14}\\
$G(\vel)$; $\alpha,\beta,\gamma$ & DWA objective and the weights of heading, clearance and velocity & \cref{ch:ch14}\\
$U(\pos)$; $\rho$, $\rho_0$ & potential field; clearance to the nearest obstacle and influence distance & \cref{ch:ch15}\\
\bottomrule
\end{tabularx}
\end{center}

\section*{Probability and filtering}

\begin{center}\small
\begin{tabularx}{\textwidth}{@{}lYl@{}}
\toprule
\textbf{Symbol} & \textbf{Meaning} & \textbf{Where introduced}\\
\midrule
$\Prob(\cdot)$, $p(\state)$ & probability of an event; probability density of a random vector & \cref{ch:ch02}\\
$\E[\cdot]$, $\Var(\cdot)$, $\Cov(\cdot)$ & expectation, variance, covariance & \cref{ch:ch02}\\
$\Normal(\vect{\mu},\mat{\Sigma})$ & Gaussian with mean $\vect{\mu}$ and covariance $\mat{\Sigma}$ & \cref{ch:ch02}\\
$\sigma_i$, $\rho_{ij}$ & standard deviation of a component; correlation coefficient & \cref{ch:ch02}\\
$\mathcal{E}_k$ & $k\sigma$ covariance ellipse, the points at Mahalanobis distance $k$ & \cref{ch:ch02}\\
$\state_k$, $\meas_k$ & state and measurement at step $k$; $\meas_{1:k}$ all measurements up to $k$ & \cref{ch:ch02}\\
$\hat{\state}_k^-$, $\mat{P}_k^-$ & predicted mean and covariance, before $\meas_k$ is used & \cref{ch:ch18}\\
$\hat{\state}_k$, $\mat{P}_k$ & updated (posterior) mean and covariance & \cref{ch:ch18}\\
$\mat{F}_k$ & state-transition matrix; for a nonlinear $f$, the Jacobian $\partial f/\partial\state$ & \cref{ch:ch18,ch:ch19}\\
$\mat{H}_k$ & measurement matrix; for a nonlinear $h$, the Jacobian $\partial h/\partial\state$ & \cref{ch:ch18,ch:ch19}\\
$\mat{Q}_k$, $\mat{R}_k$ & process-noise and measurement-noise covariance & \cref{ch:ch18}\\
$\vect{w}_k$, $\vect{v}_k$ & process noise $\Normal(\vect{0},\mat{Q}_k)$ and measurement noise $\Normal(\vect{0},\mat{R}_k)$ & \cref{ch:ch19}\\
$\vect{y}_k$, $\mat{S}_k$ & innovation $\meas_k-h(\hat{\state}_k^-)$ and its covariance $\mat{H}_k\mat{P}_k^-\mat{H}_k\T+\mat{R}_k$ & \cref{ch:ch19}\\
$\Kgain_k$ & Kalman gain $\mat{P}_k^-\mat{H}_k\T\mat{S}_k\inv$ & \cref{ch:ch18,ch:ch19}\\
$\mathcal{X}_i$, $W_i^{(m)}$, $W_i^{(c)}$ & sigma points of the unscented transform and their mean and covariance weights & \cref{ch:ch19}\\
$\state^{(i)}$, $w^{(i)}$ & particle $i$ and its weight; $N_{\mathrm{eff}}=1/\sum_i (w^{(i)})^2$ & \cref{ch:ch19}\\
$T_{\mathrm{obs}}$, $T_{\mathrm{pred}}$ & length of the observation window and of the prediction horizon & \cref{ch:ch20}\\
$h$, $\hat{\pos}_{T+h}$ & horizon step $h\in\set{1,\dots,T_{\mathrm{pred}}}$ and the position predicted for it & \cref{ch:ch20}\\
$\mathrm{ADE}$, $\mathrm{FDE}$ & average and final displacement error; $\mathrm{minADE}_K$ over $K$ samples & \cref{ch:ch20}\\
$\mat{Q}$, $\mat{K}$, $\mat{V}$ & query, key and value matrices of attention (this meaning only in \cref{ch:ch20}) & \cref{ch:ch20}\\
\bottomrule
\end{tabularx}
\end{center}

\paragraph{Two warnings about letters.}
$\mat{Q}$ and $\mat{R}$ are \emph{noise covariances} in Part~VI, where they say
how much the motion model and the sensor are trusted, and \emph{cost weights} in
Part~VII, where they say how much a tracking error and a control effort cost.
The two roles are unrelated, they have different units, and no equation in this
book uses both at once; \cref{ch:ch21} repeats the warning where the weights are
introduced. \Cref{ch:ch20} adds a third, local, meaning: inside the section on
attention, $\mat{Q}$, $\mat{K}$ and $\mat{V}$ are the query, key and value
matrices. Second, this book follows the estimation literature and writes
$\mat{F}$ for the state-transition matrix and $\mat{H}$ for the measurement
matrix, while control texts usually write $\mat{A}$ and $\mat{C}$ for the same
two objects; the motion models of \cref{ch:ch02} and the controller of
\cref{ch:ch21} do use $\mat{A}$ and $\mat{B}$, because there they are the model
of a plant and not of a filter.

\section*{Optimization and control}

\begin{center}\small
\begin{tabularx}{\textwidth}{@{}lYl@{}}
\toprule
\textbf{Symbol} & \textbf{Meaning} & \textbf{Where introduced}\\
\midrule
$J$, $J^*$ & objective (cost) of an optimization problem and its optimal value & \cref{ch:ch21,ch:ch22}\\
$\st$ & subject to: the constraints of an optimization problem follow & \cref{ch:ch21}\\
$N$ & prediction horizon of the controller, in steps of $\dt$ & \cref{ch:ch21}\\
$\state_k^{\mathrm{ref}}$ & reference state at step $k$ of the horizon ($k=0$ is now) & \cref{ch:ch21}\\
$\mat{Q}\succeq\mat{0}$, $\mat{R}\succ\mat{0}$ & weights on the state error and on the control effort & \cref{ch:ch21}\\
$\mat{P}\succeq\mat{0}$ & terminal cost weight on the last predicted state & \cref{ch:ch21}\\
$\mat{S}_x$, $\mat{S}_u$ & prediction matrices that stack the horizon into one linear map & \cref{ch:ch21}\\
$\vect{z}$ & stacked decision variable of the resulting quadratic or linear program & \cref{ch:ch21,ch:ch22}\\
$\mat{H}$, $\vect{f}$ & Hessian and linear term of the quadratic program & \cref{ch:ch21}\\
$\mat{G}$, $\vect{l}$, $\vect{u}$ & constraint matrix and its lower and upper bounds, $\vect{l}\le\mat{G}\vect{z}\le\vect{u}$ & \cref{ch:ch21}\\
$\vect{y}$, $\rho$ & Lagrange multipliers and the penalty parameter of the ADMM solver & \cref{ch:ch21}\\
$s_k$ & slack variable that softens a constraint at step $k$ & \cref{ch:ch21}\\
$\eps$ & tolerance of a numerical solver (integrality, primal and dual residuals) & \cref{ch:ch21,ch:ch22}\\
$\vect{c}\T\vect{z}$ & objective of a linear program; $\mat{A}\vect{z}\le\vect{b}$ its inequalities & \cref{ch:ch22}\\
$P$ & feasible set of a linear program, a polyhedron & \cref{ch:ch22}\\
$z_j\in\Z$, $b\in\set{0,1}$ & integer variable and binary variable of a mixed-integer program & \cref{ch:ch22}\\
$M$ & big-$M$ constant that switches an inequality off when its binary is $1$ & \cref{ch:ch22}\\
\bottomrule
\end{tabularx}
\end{center}

\section*{Networks and consensus}

\begin{center}\small
\begin{tabularx}{\textwidth}{@{}lYl@{}}
\toprule
\textbf{Symbol} & \textbf{Meaning} & \textbf{Where introduced}\\
\midrule
$G=(V,E)$ & communication graph: the drones and the pairs that can exchange messages & \cref{ch:ch23}\\
$N_i$ & neighbours of drone $i$ in the communication graph & \cref{ch:ch23}\\
$\mat{A}=(a_{ij})$ & adjacency matrix of the communication graph & \cref{ch:ch23}\\
$d_i$, $\mat{D}$ & degree of drone $i$; degree matrix $\diag(d_1,\dots,d_n)$, largest degree $d_{\max}$ & \cref{ch:ch23}\\
$\mat{L}=\mat{D}-\mat{A}$ & graph Laplacian; symmetric, positive semidefinite, $\mat{L}\vect{1}=\vect{0}$ & \cref{ch:ch23}\\
$\lambda_2$, $\lambda_n$ & algebraic connectivity (Fiedler value) and largest Laplacian eigenvalue & \cref{ch:ch23}\\
$\bar{x}$, $\vect{\delta}$ & average of the states; disagreement vector $\vect{x}-\bar{x}\vect{1}$ & \cref{ch:ch23}\\
$\eps$ & step size of discrete-time consensus, $\vect{x}_{k+1}=(\mat{I}-\eps\mat{L})\vect{x}_k$ & \cref{ch:ch23}\\
$R_{\mathrm{comm}}$ & communication range: $\set{i,j}\in E$ exactly when $\norm{\pos_i-\pos_j}\le R_{\mathrm{comm}}$ & \cref{ch:ch21,ch:ch23}\\
$\vect{o}_i$, $\vect{d}_{ij}$ & offset of drone $i$ in the formation; desired displacement $\vect{o}_j-\vect{o}_i$ & \cref{ch:ch23}\\
$e_F$, $e_{\max}$ & formation error (RMS violation of the displacements) and its tolerance & \cref{ch:ch21,ch:ch23}\\
$k_p$, $k_v$, $k_d$ & position, velocity and damping gains of second-order consensus & \cref{ch:ch23}\\
\bottomrule
\end{tabularx}
\end{center}
